
我们将引入\textbf{\trSymplecticManifold}（symplectic manifold）、%
\textbf{\trSymplecticMapping}（symplectic mapping）与%
\textbf{\trHamiltonianVectorField}（Hamiltonian vector field）这几个概念。%
我们无意对\textbf{\trHamiltonianFormalism}（Hamiltonian formalism）作系统的处理，因为它已在许多书中得到阐述。%
毋宁说，我们将围绕这些概念提出若干问题，它们近来引出了新的现象与有趣的未解决问题。%
例如，“一个\trSymplecticMapping{}能做些什么？”这一问题，便引出了不同于\textbf{\trVolume}（volume）的新的%
\textbf{\trSymplecticInvariant}（symplectic invariant），%
这些不变量将在后续各章中得到详细讨论。%
我们将说明，一个起源于\textbf{\trCelestialMechanics}（celestial mechanics）、看似截然不同却历史悠久的问题，%
其实也与这些不变量相关。具体而言，受\trPoincare{}\trRecurrence{}定理的启发，我们要问：%
\trHamiltonianVectorField{}的\textbf{\trCompact\trEnergySurface}（compact energy surface）是否具有\textbf{\trPeriodicOrbit}（periodic orbit）。%
对于 $\mathbb{R}^{2n}$ 中的\textbf{\trConvex\trHypersurface}（convex hypersurface）这一非常特殊的情形——%
它在\trHamiltonian{}系统的这一定性问题历史上堪称里程碑之一——%
我们将给出一个存在性证明，以阐明\textbf{\trCalculusOfVariations}（calculus of variations）中所谓的%
\textbf{\trDirectMethod}（direct method）。%
这一经典方法，与后续各章中为建立整体\textbf{\trPeriodicSolution}（periodic solution）而引入的较新方法，恰成对照。%
在本章的末尾，我们将在不给证明的情形下展示，%
一个给定的\trPeriodicOrbit{}附近可以预期的丰富而错综复杂的\textbf{\trOrbit}（orbit）结构。%
这些考察基于详加介绍的局部\trNonlinear\trBirkhoffInvariant。
